
Soit donc @@M0@@ un torseur sous @@M1@@, correspondant à une extension @@M2@@ de @@M3@@ par @@M4@@. L'algèbre @@M5@@ de @@M6@@ sur @@M7@@ est @@M8@@, soit @@M9@@, qui s'identifie à l'anneau des fractions homogènes de degré @@M10@@ @@M11@@. Par suite, @@M12@@ s'identifie à l'ouvert @@M13@@ de @@M14@@. Le complémentaire de @@M15@@ dans @@M16@@ est alors @@M17@@, qui s'identifie d'après (3.3.1) à @@M18@@. On a donc les immersions :
@@M19@@
où @@M20@@ est fermé dans @@M21@@, @@M22@@ ouvert et complémentaire de @@M23@@.

\textbf{Proposition 3.4.} \uline{Soient} @@M24@@ \uline{un schéma noethérien}, @@M25@@ \uline{un} @@M26@@\uline{-Module localement libre de rang} @@M27@@, @@M28@@ \uline{le fibré vectoriel associé}, @@M29@@ \uline{un torseur sous} @@M30@@, @@M31@@ \uline{le morphisme structural de} @@M32@@. \uline{Alors l'homomorphisme canonique} :
@@M33@@
\uline{est un isomorphisme}.

Avec les identifications de 3.3, on a la suite exacte :
@@M34@@
où @@M35@@ désigne l'immersion fermée @@M36@@ définie par @@M37@@. D'après 3.2, on a @@M38@@, @@M39@@.
Comme @@M40@@ et @@M41@@ sont lisses, @@M42@@ est une immersion régulière, et on a donc un homomorphisme @@M43@@; l'homomorphisme entre les @@M44@@ est alors obtenu en tensorisant l'homomorphisme entre les @@M45@@ par @@M46@@ : il suffit pour le voir de montrer que pour @@M47@@ et
